\section{Strategies}

\subsection{Maslow's Hierarchy of Needs: The Gradient and Dynamic of Needs}

The ``Hierarchy of Needs'' proposed by Abraham Maslow in 1943 is one of
the most influential models in psychology history.
He depicted human needs as a progressive system from basic survival to
advanced spiritual realization:

\begin{itemize}
    \item \textbf{Physiological needs:} Breathing, food, sleep. This is
    the bottom cornerstone of survival \citep{damasio1994}.
    \item \textbf{Safety needs:} Personal security, health guarantee,
    resource ownership \citep{goleman1995}.
    \item \textbf{Love and belonging:} Friendship, love, intimacy.
    This is the core connection most frequently touched upon in communication \citep{ainsworth1978}.
    \item \textbf{Esteem:} Self-respect, sense of achievement, recognition
    by others \citep{erikson1950}.
    \item \textbf{Self-actualization:} Creativity, morality, problem-solving
    ability, unleashing one's potential.
\end{itemize}

\paragraph{Deficit Needs vs. Being Needs}
Maslow collectively referred to the first four levels as \textbf{``Deficit
Needs.''}
They are like the body's ``hunger''; once missing, the individual falls
into tension and anxiety, and all attention is drawn to fill these voids.
And self-actualization is a \textbf{``Being Need,''} it does not originate
from deficiency, but from the desire for life's abundance \citep{salovey1990}.

\paragraph{A Dynamic Flowing Perspective}
Although the pyramid structure seems like steps, Maslow emphasized in his
later years that the satisfaction of needs is \textbf{dynamic and highly
fluid.}
A painter who is willing to endure hunger for artistic creation
(self-actualization) proves that higher-level needs can become the driving
force for lower-level needs \citep{hoffman2000}.
In relationship conflicts, understanding this is crucial:
When the other person shows anger or attack, he is often due to ``esteem''
or ``belonging'' deficit needs being threatened, retreating into a defensive
state.
\textbf{Identifying which level of the pyramid the other person is currently
standing on is the key to providing ``precise empathy.''}

\subsection{Self-Determination Theory}
The Self-Determination Theory (SDT) proposed by Edward Deci and Richard
Ryan argues that three core psychological needs are the basis of human
intrinsic motivation:
\begin{itemize}
    \item \textbf{Autonomy:} Feeling that one's behavior is out of free
    choice, rather than forced obedience \citep{thomas1977}.
    \item \textbf{Competence:} Feeling capable of dealing with challenges
    in the environment and operating effectively \citep{dweck2006}.
    \item \textbf{Relatedness:} Feeling authentically connected with others,
    feeling cared for and caring for others \citep{feldman2024}.
\end{itemize}
Studies have shown that when these three major needs are consistently met
in a relationship, individuals will show higher intrinsic motivation,
better mental health status, and more stable relationship satisfaction.
Conversely, the long-term lack of any one need will lead to psychological
sluggishness, relationship conflict, and inner emptiness.

\begin{table}[h]
\centering
\caption{Universal Human Needs}
\label{tab:universal_needs}
\begin{tabular}{|l|p{11cm}|}
\hline
\textbf{Category} & \textbf{Core Needs} \\ \hline
Connection & Acceptance, belonging, intimacy, love, companionship, communication, understanding, contribution, mutuality, security, celebration \\ \hline
Physical Well-Being & Air, food, water, shelter, touch, rest, movement, sex, safety (physiological), health \\ \hline
Integrity & Authenticity, honesty, consistency \\ \hline
Play & Fun, humor, ease, joy \\ \hline
Peace & Beauty, harmony, inspiration, order, tranquility \\ \hline
Autonomy & Choice, freedom, independence, space, spontaneity \\ \hline
Meaning & Awareness, challenge, clarity, competence, creativity, efficacy, growth, hope, learning, mourning, purpose, self-expression \\ \hline
\end{tabular}
\end{table}

As shown in Table \ref{tab:universal_needs}, human needs are multi-layered
and multi-dimensional.
When we can replace judgments with the identification of needs in daily
communication, many seemingly irreconcilable conflicts will show new
possibilities.

\begin{keyinsight}
Human needs range from basic survival and deficit needs to profound desires
for autonomy, competence, and relatedness.
By recognizing that these fundamental psychological drives shape our internal
motivations and relationship dynamics, we can precisely identify the roots of
conflict and respond with targeted empathy rather than reactive judgment.
\end{keyinsight}
